\chapter{Conclusion and Future Work}
\label{chap:conclusion}

\section{Summary of the Work}

AURA-MAS separates surveillance into a sequence of identifiable steps. Camera
and audio agents emit structured events. The fusion agent groups related events
into hypotheses, the coordinator may request verification, and the policy agent
makes the alert decision. Explanations are generated only after that decision.
This arrangement keeps raw video at the edge and prevents generative output
from changing the alert path.

The final evaluation used 373 replay runs across nine scenario manifests. It
also kept the pre-fix baseline and documented every change between the two
measurement pipelines in Appendix~\ref{app:methodology}. The results do not
support a performance ranking of the four architectures. After paired Wilcoxon
tests and Holm--Bonferroni correction, none of the six pairwise F1 differences
is statistically significant.

RQ1 remains unresolved. The 80.4\% reduction in time-to-alert occurred only
in \texttt{demo\_site\_01} ($n=5$, $p=0.0625$). Across the seven scenarios
with all four modes, the reduction was 10.7\% ($p=0.59$). The replay driver
also paces MAS modes but not the centralized mode, so the comparison mixes
architecture with wall-clock behaviour of the fusion window and cooldown.

RQ2 cannot validate the auction as a spatial allocation method. The campaign
manifests did not supply field-of-view overlap, so the bid's informative term
always used its 0.5 fallback. In a two-camera setting, that does not test
whether the selected camera has the best viewpoint. The no-coordination and
auction comparison is not significant after correction.

RQ3 confirms that cross-modal fusion can occur. Four of 240 campaign alerts
include both a camera and a microphone. Its effect on confidence cannot be
estimated from this corpus, however. The often-cited 0.78 alert was a
single-camera video alert, and repeated events from one camera make the
noisy-OR score reach one before an audio contribution can be seen.

RQ4 is supported as a design property, not as a campaign result. Unit tests
and one post-campaign run show that the guardrail can reject an unsupported
report. The campaign itself never enabled \texttt{--llm}, so all 373
explanations came from the deterministic template. RQ5 has the clearest
evidence: the campaign wrote 1,737 anonymized crops and 373 audit files. The
audit log is append-only, not cryptographically immutable, and the usefulness
of the alerts for operators remains unmeasured.

\section{Contributions Revisited}

The final evaluation narrows the claims for the five contributions announced
in Chapter~\ref{chap:introduction}.

\begin{description}
    \item[C1 --- Architecture.] The system implements six agents and an
    operator console that separate perception, coordination, policy,
    explanation, and oversight. MQTT and Redis are available as transports,
    but the campaign used \texttt{LocalBus}, so it did not measure network
    cost.

    \item[C2 --- Active verification auction.] The task, bid, award, and
    verification flow is implemented and testable. Because campaign manifests
    left the spatial utility term unpopulated, the reported results do not show
    an advantage over round-robin allocation.

    \item[C3 --- Late fusion.] The noisy-OR implementation is monotone and
    accepts multi-sensor input. At the observed operating confidences it can
    saturate, and repeated events from one camera violate the model's
    independence assumption.

    \item[C4 --- Guarded explanation.] The finite-state workflow and evidence
    checks keep explanations separate from alert decisions. The campaign did
    not exercise the LLM path, so it provides no evidence about explanation
    quality or groundedness in normal operation.

    \item[C5 --- Evaluation apparatus.] Replay manifests, isolated
    subprocesses, repetition plumbing, serialized audit artifacts, the
    preserved baseline, and the methodology-change log make the results
    reproducible, including the negative findings.
\end{description}

\section{Limitations}

The paired analysis has only about nine scenario-level observations. The
observed no-coordination versus auction effect ($d_z=0.347$) would need about
65 paired observations for 80\% power, or roughly 13 scenarios at the current
repetition design. More than twenty implementation parameters influence the
results, but the project has no full held-out tuning protocol or seed control.

The metrics also set limits on interpretation. Event-family matching with a
$\pm5$ s tolerance is a system-level measure, not a per-frame detection or
tracking measure. The corpus has no per-frame track ground truth, so HOTA,
MOTA, and IDF1 cannot be calculated. False alerts per hour are unstable for
runs that last only seconds.

This remains a demonstrator-scale evaluation. Only \texttt{demo\_site\_01}
is a multi-zone manifest from the campaign; later exploratory multi-zone
manifests are outside the 373-run evidence. The CLIP anomaly scorer has an AUC
of 0.308 in its calibration note. A post-campaign appearance-descriptor
prototype was withheld from the evaluated configuration. It produced no
recorded Re-ID matches and did not persist Re-ID features, so the thesis keeps
the stated exclusion of person re-identification.

\section{Future Work}

The auction requires field-of-view overlap derived from camera calibration or
homography, followed by a comparison with a paced centralized control and
stronger allocation baselines such as Hungarian assignment. Verification
should examine evidence from the same time as the hypothesis. Fusion should
pool dependent observations before combining them.

The scenario corpus should grow to at least thirteen independent scenarios,
with controlled seeds, held-out parameter selection, and richer annotations.
Per-frame track ground truth would permit HOTA, MOTA, and IDF1. A physical
deployment could then measure MQTT/Redis transport cost, edge resources, and
operator response instead of relying on in-process replay.

The explanation pipeline needs a campaign with \texttt{--llm} enabled, the
judge scaffolding executed, and an operator study that separates groundedness
from usefulness. A hash chain or signed audit store would make the present
append-only log tamper-evident.

\section{Closing Remarks}

The campaign did not confirm the strongest claims from the first
demonstration. It did leave a traceable record of why they failed, what was
measured, and what should be tested next. That record is necessary before the
prototype can support stronger claims.
